\enquestion{9.42}{
    A melting point test of $n = 15$ samples of a binder used in manufacturing a rocket propellant resulted in $\bar{x} = \qty{68}{\degree C}$.
    Assume that the melting point is normally distributed with $\sigma = \qty{0.83}{\degree C}$.
    \begin{enumerate}[label=(\alph*)]
        \item Test $H_{0}: \mu = 68.3$ versus $H_{1}: \mu \ne 68.3$ using $\alpha = 0.01$.
        \item What is the $p$-value for this test?
        \item What is the $\beta$-error if the true mean is $\mu = 65.5$?
    \end{enumerate}
}

\enquestion{9.47}{
    Medical researchers have developed a new artificial heart constructed primarily of titanium and plastic.
    The heart will last and operate almost indefinitely once it is implanted in the patient's body, but the battery pack needs to be recharged about every four hours.
    A random sample of $50$ battery packs is selected and subjected to a life test.
    The average life of these batteries is $\qty{4.05}{\hour}$.
    Assume that battery life is normally distributed with standard deviation $\sigma = \qty{0.2}{\hour}$.
    \begin{enumerate}[label=(\alph*)]
        \item Is there evidence to support the claim that mean battery life exceeds $\qty{4}{\hour}$?
        Use $\alpha = 0.05$.
        \item What is the $p$-value for the test in part (a)?
        \item Compute the power of the test if the true mean battery life is $\qty{4.5}{\hour}$.
    \end{enumerate}
}

\enquestion{9.51}{
    For the hypothesis test $H_{0}: \mu = 7$ against $H_{1}: \mu \ne 7$ with variance unknown and $n = 25$, approximate the $p$-value for each of the following test statistics.
    
    \begin{center}
        \begin{tabular}{*{3}{C{0.3\textwidth}}}
            (a) $t_{0} = 2.05$ & (b) $t_{0} = -1.84$ & (c) $t_{0} = 0.4$
        \end{tabular}
    \end{center}
}

\enquestion{9.59}{
    A 1992 article in the \textit{Journal of the American Medical Association} (\enquote{A Critical Appraisal of 98.6 Degrees F, the Upper Limit of the Normal Body Temperature, and Other Legacies of Carl Reinhold August Wunderlich}) reported body temperature, gender, and heart rate for a number of subjects.
    The body temperatures for $25$ female subjects follow: $97.8$, $97.2$, $97.4$, $97.6$, $97.8$, $97.9$, $98.0$, $98.0$, $98.0$, $98.1$, $98.2$, $98.3$, $98.3$, $98.4$, $98.4$, $98.4$, $98.5$, $98.6$, $98.6$, $98.7$, $98.8$, $98.8$, $98.9$, $98.9$, and $99.0$.
    \begin{enumerate}
        \item[(a)] Test the hypothesis $H_{0}: \mu = 98.6$ versus $H_{1}: \mu \ne 98.6$, using $\alpha = 0.05$.
        Find the $p$-value.
        \item[(e)] Explain how the question in part (a) could be answered by constructing a two-sided confidence interval on the mean female body temperature.
    \end{enumerate}
}

\enquestion{9.60}{
    Cloud seeding has been studied for many decades as a weather modification procedure (for an interesting study of this subject, see the article in \textit{Technometrics}, \enquote{A Bayesian Analysis of a Multiplicative Treatment Effect in Weather Modification,} Vol. 17, pp. 161–166).
    The rainfall in cubic-meters from $20$ clouds that were selected at random and seeded with silver nitrate follows: $22.2 \times 10^{3}$, $37.9 \times 10^{3}$, $24.4 \times 10^{3}$, $33.4 \times 10^{3}$, $27.5 \times 10^{3}$, $23.2 \times 10^{3}$, $39.2 \times 10^{3}$, $28.9 \times 10^{3}$, $26.1 \times 10^{3}$, $34.4 \times 10^{3}$, $39.3 \times 10^{3}$, $33.4 \times 10^{3}$, $30.8 \times 10^{3}$, $30.4 \times 10^{3}$, $33.2 \times 10^{3}$, $26.9 \times 10^{3}$, $36.2 \times 10^{3}$, $42.9 \times 10^{3}$, $32.9 \times 10^{3}$, and $38.9 \times 10^{3}$.
    \begin{enumerate}
        \item[(a)] Can you support a claim that the mean rainfall from seeded clouds exceeds $\qty{30e3}{\m\cubed}$?
        Use $\alpha = 0.01$.
        Find the $p$-value.
        \item[(e)] Explain how the question in part (a) could be answered by constructing a one-sided confidence bound on the mean diameter.
    \end{enumerate}
}

\enquestion{9.126}{
    Consider the following computer output.

    \begin{center}
        \begin{tabular}{*{7}{C{0.11428571428\textwidth}}}
            \toprule
                \multicolumn{7}{l}{\textbf{One-Sample $Z$:}} \\                
                \multicolumn{7}{l}{Test of $\mu = 26$ vs $\mu>26$} \\
                \multicolumn{7}{l}{The assumed standard deviation $= 1.7$} \\   
                Variable & $n$ & Mean & StDev & SE Mean & $Z$ & $p$ \\
                $X$ & ? & $26.541$ & $2.032$ & $0.401$ & ? & ? \\                   
            \bottomrule
        \end{tabular}
    \end{center}

    \begin{enumerate}[label=(\alph*)]
        \item Fill in the missing information.
        \item Is this a one-sided or a two-sided test?
        \item What are your conclusions if $\alpha = 0.05$?
        \item Find a \SI{95}{\percent} two-sided CI on the mean.
    \end{enumerate}
}

\enquestion{9.127}{
    Consider the following computer output.

    \begin{center}
        \begin{tabular}{*{8}{C{0.1\textwidth}}}
            \toprule
                \multicolumn{8}{l}{\textbf{One-Sample $T$:}} \\                
                \multicolumn{8}{l}{Test of $\mu = 100$ vs $\mu\neq 100$} \\
                Variable & $n$ & Mean & StDev & SE Mean & \SI{95}{\percent} CI & $t$ & $p$ \\
                $X$ & $16$ & $98.33$ & $4.61$ & ? & $($?,?$)$ & ? & ? \\                   
            \bottomrule
        \end{tabular}
    \end{center}
    
    \begin{enumerate}[label=(\alph*)]
        \item How many degrees of freedom are there on the $t$-statistic?
        \item Fill in the missing information.
        You may use bounds on the $p$-value.
        \item What are your conclusions if $\alpha = 0.01$?
        \item What are your conclusions if the hypothesis is $H_{0}: \mu = 100$ versus $H_{0}: \mu > 100$?
    \end{enumerate}
}